% !TeX program = latexmk
% !BIB program = biber
% !TeX spellcheck = ru_RU
\documentclass[14pt]{extarticle}
\usepackage{iftex}
\ifPDFTeX
  \usepackage[T2A]{fontenc}
  \usepackage[utf8]{inputenc}
\else
  \usepackage{fontspec}
  \setmainfont{DejaVu Serif}
\fi
\usepackage[left=2cm,right=2cm, top=2cm,bottom=2cm]{geometry}
\usepackage[english,russian]{babel}
\usepackage[unicode,colorlinks,filecolor=black,citecolor=black,linkcolor=black]{hyperref}
\usepackage{booktabs}
\usepackage{amsmath} 
\usepackage{amssymb}
\usepackage{listings}
\usepackage{xcolor}
\usepackage{graphicx}
\graphicspath{ {./} }
\usepackage{enumitem}
\usepackage{siunitx}
\usepackage[font={footnotesize,it}]{caption}
\usepackage{tabularx}
\usepackage{float}
\usepackage[backend=biber,style=numeric,sorting=none]{biblatex}
\addbibresource{references.bib}
\usepackage{titlesec}

\titleformat{\section}
  {\normalfont\large\bfseries\raggedright}
  {\thesection.}                   
  {1em}                             
  {}
\titleformat{\subsection}
  {\normalfont\normalsize\itshape\bfseries}  
  {\thesubsection.}                 
  {1em}                            
  {}
  \titleformat{\subsubsection}
  {\normalfont\small\bfseries}  
  {\thesubsubsection.}               
  {1em}                            
  {}           

\definecolor{codegreen}{rgb}{0,0.6,0}
\definecolor{codeblue}{rgb}{0,0,0.6} 
\definecolor{codepurple}{rgb}{0.58,0,0.82} 
\definecolor{backcolour}{rgb}{0.98,0.98,0.98} 

\lstdefinestyle{json}{
    backgroundcolor=\color{backcolour},   
    commentstyle=\color{codegreen},
    stringstyle=\color{codepurple},
    numberstyle=\color{codeblue}\footnotesize,
    basicstyle=\ttfamily\footnotesize,
    breakatwhitespace=false,         
    breaklines=true,                 
    captionpos=b,                    
    keepspaces=true,                 
    showspaces=false,                
    showstringspaces=false,
    showtabs=false,                  
    tabsize=2,
    frame=single,                    
    rulecolor=\color{black}
}
\setcounter{secnumdepth}{5}
\setcounter{biburllcpenalty}{7000}
\setcounter{biburlucpenalty}{8000}
\setlength{\parskip}{10pt}
\begin{document}
\sloppy
\begin{titlepage}
\begin{center}
\normalsize
\normalsize{ФЕДЕРАЛЬНОЕ ГОСУДАРСТВЕННОЕ АВТОНОМНОЕ\\ОБРАЗОВАТЕЛЬНОЕ УЧРЕЖДЕНИЕ\\ВЫСШЕГО ОБРАЗОВАНИЯ\\«НАЦИОНАЛЬНЫЙ ИССЛЕДОВАТЕЛЬСКИЙ УНИВЕРСИТЕТ\\

ВЫСШАЯ ШКОЛА ЭКОНОМИКИ»}

\vfill

\textbf{Факультет информатики, математики и компьютерных наук}\\[3mm]

\textbf{Программа подготовки бакалавров по направлению\\Компьютерные науки и технологии}

\vfill

\textit{Пестов Лев Евгеньевич}\\[3mm]

\textbf{КУРСОВАЯ РАБОТА}\\[10mm]

\normalsize{ 
 	
Разработка системы семантической сегментации дорожной сцены с акцентом на сегментацию слабоконтрастных объектов}

\end{center}

\vfill
\newlength{\ML}
\settowidth{\ML}{«\underline{\hspace{0.7cm}}»
\underline{\hspace{2cm}}}
\hfill
\begin{minipage}{0.4\textwidth}
\raggedleft{Научный руководитель\\
преподаватель НИУ ВШЭ - НН\\[2mm]
Сенников Андрей Олегович}
\end{minipage}
\vfill
\begin{center}

Нижний Новгород, 2026г.

\end{center}
\end{titlepage}
\newpage  

\renewcommand{\contentsname}{Структура работы}
\tableofcontents
\newpage

\section{Введение}

\textbf{Актуальность темы.} Развитие систем автономного вождения и помощи водителю (ADAS) требует высокой надежности алгоритмов компьютерного зрения в любых условиях эксплуатации. Современные методы семантической сегментации демонстрируют высокие показатели качества (mIoU > 80\%) на стандартных бенчмарках, таких как Cityscapes, содержащих изображения, полученные в благоприятных погодных условиях и при хорошем освещении. Однако реальная дорожная обстановка часто сопряжена с факторами, снижающими визуальное качество сцены: туман, дождь, снегопад, сумерки или ночное время суток. В таких условиях контрастность объектов критически снижается, что приводит к значительной деградации качества работы нейронных сетей, обученных на «чистых» данных.

Особую опасность представляет потеря детектирования слабоконтрастных объектов — пешеходов в темной одежде, велосипедистов или удаленных транспортных средств, сливающихся с фоном. Ошибка сегментации в таких случаях может привести к аварийным ситуациям. При этом сбор и разметка масштабных датасетов для всех возможных погодных условий является трудоемкой и дорогостоящей задачей. В связи с этим актуальной становится задача разработки методов адаптации моделей, обученных на качественных данных, к сложным условиям видимости без необходимости полной переразметки данных.

\textbf{Проблема исследования} заключается в существенном падении точности (Domain Shift) современных алгоритмов семантической сегментации при переходе от домена с хорошей видимостью к домену с неблагоприятными погодными условиями, особенно в части распознавания малых и слабоконтрастных объектов.

\textbf{Цель работы} — разработка и исследование методики повышения качества семантической сегментации дорожных сцен в сложных погодных условиях (туман, ночь, дождь) с использованием методов доменной адаптации и модификации процесса обучения.

Для достижения поставленной цели необходимо решить следующие \textbf{задачи}:
\begin{enumerate}
    \item Провести обзор современных архитектур нейронных сетей (CNN и Transformer-based) для задачи семантической сегментации.
    \item Изучить проблему доменного сдвига (Domain Shift) и методы повышения обобщающей способности моделей (Data Augmentation, Domain Adaptation).
    \item Реализовать базовый пайплайн обучения (baseline) на наборе данных Cityscapes и оценить его качество на данных со сложными погодными условиями (ACDC).
    \item Разработать и внедрить модификации (аугментации, функции потерь), направленные на улучшение распознавания слабоконтрастных объектов.
    \item Провести сравнительный анализ результатов и оценить эффективность предложенных методов.
\end{enumerate}

\textbf{Объект исследования:} алгоритмы семантической сегментации изображений дорожных сцен.

\textbf{Предмет исследования:} методы повышения устойчивости нейронных сетей к изменению визуальных условий (освещение, погода) и снижению контраста.

\newpage

\section{Теоретические основы семантической сегментации и доменной устойчивости}

\subsection{Семантическая сегментация дорожной сцены}

\subsubsection{Формализация задачи pixel-wise классификации}
Семантическая сегментация представляет собой задачу плотной (pixel-wise) классификации, в которой каждому пикселю входного изображения сопоставляется метка класса из заранее заданного множества. Формально, пусть задано входное изображение $X \in \mathbb{R}^{H \times W \times C}$, где $H$ и $W$ --- пространственные размеры изображения (высота и ширина), а $C$ --- количество каналов (обычно $C=3$ для RGB-изображений). Пусть $\mathcal{K} = \{1, 2, \dots, K\}$ --- множество целевых семантических классов, характерных для дорожной сцены (например, дорога, тротуар, здания, транспорт, пешеходы и т.д.).

Цель модели семантической сегментации (нейронной сети), параметризованной весами $\theta$, заключается в поиске отображения $f_\theta$:
\begin{equation}
    f_\theta : \mathbb{R}^{H \times W \times C} \to \mathbb{R}^{H \times W \times K},
\end{equation}
которое преобразует входное изображение в тензор логитов $Z \in \mathbb{R}^{H \times W \times K}$. Для каждого пространственного положения (пикселя) с координатами $(i, j)$ сеть предсказывает вектор $z_{i,j} \in \mathbb{R}^K$.

Вероятностное распределение принадлежности пикселя к каждому из $K$ классов получается путем применения функции \textit{softmax} вдоль канала классов:
\begin{equation}
    p_{i,j}^{(k)} = \frac{\exp\left(z_{i,j}^{(k)}\right)}{\sum_{c=1}^K \exp\left(z_{i,j}^{(c)}\right)},
\end{equation}
где $p_{i,j}^{(k)}$ --- предсказанная вероятность того, что пиксель $(i, j)$ принадлежит классу $k$.

Итоговая предсказанная маска сегментации $\hat{Y} \in \mathcal{K}^{H \times W}$ формируется путем выбора класса с максимальной вероятностью для каждого пикселя:
\begin{equation}
    \hat{y}_{i,j} = \arg\max_{k \in \mathcal{K}} p_{i,j}^{(k)}.
\end{equation}

В процессе обучения модели наиболее распространенным функционалом потерь является средняя категориальная кросс-энтропия (Categorical Cross-Entropy, CCE). Если $y_{i,j}$ --- истинная метка класса для пикселя $(i, j)$, представленная в виде one-hot вектора (где $y_{i,j}^{(k)} = 1$, если пиксель принадлежит классу $k$, и $0$ иначе), то функция потерь имеет вид:
\begin{equation}
    \mathcal{L}_{CE} = - \frac{1}{H \cdot W} \sum_{i=1}^{H} \sum_{j=1}^{W} \sum_{k=1}^{K} y_{i,j}^{(k)} \log\left(p_{i,j}^{(k)}\right).
\end{equation}

В отличие от задач классификации изображений в целом или детектирования объектов, сегментация требует одновременно высокой семантической точности (понимания контекста сцены) и точной локализации пространственных границ объектов. Для задач автономного вождения особенно важна стабильность качества предсказаний: модель должна корректно выделять как крупные структуры (дорожное полотно, небо), так и мелкие, слабоконтрастные, но критически важные объекты (пешеходы вдали, дорожные знаки). Это делает задачу крайне чувствительной к доменному сдвигу (domain shift) между распределением данных, на которых модель обучалась (source domain), и распределением данных в реальных условиях применения (target domain), например, при изменении погодных условий или времени суток.

\subsubsection{Основные метрики: mIoU, class-wise IoU, Pixel Accuracy}
Для количественной оценки качества моделей семантической сегментации используются метрики, основанные на вычислении площадей перекрытия предсказанных и истинных масок. В основе вычислений лежат базовые элементы матрицы ошибок (Confusion Matrix), рассчитываемые независимо для каждого класса $k$:
\begin{itemize}
    \item $TP_k$ (True Positives) --- количество истинно-положительных пикселей (пиксели, верно отнесенные к классу $k$);
    \item $FP_k$ (False Positives) --- количество ложно-положительных пикселей (пиксели, ошибочно отнесенные к классу $k$);
    \item $FN_k$ (False Negatives) --- количество ложно-отрицательных пикселей (пиксели класса $k$, отнесенные моделью к другому классу);
    \item $TN_k$ (True Negatives) --- количество истинно-отрицательных пикселей (пиксели других классов, корректно не отнесенные к классу $k$).
\end{itemize}

Ключевой метрикой качества в данной работе является \textbf{mIoU} (mean Intersection over Union --- среднее пересечение над объединением). Сначала вычисляется метрика \textbf{class-wise IoU} для каждого отдельного класса $k$, которая определяется как отношение площади пересечения предсказанной маски и ground truth маски к площади их объединения:
\begin{equation}
    \text{IoU}_k = \frac{TP_k}{TP_k + FP_k + FN_k}.
\end{equation}

Значение $\text{IoU}_k$ строго лежит в диапазоне от $0$ до $1$. Анализ class-wise IoU позволяет проводить детальную диагностику поведения модели: например, выявлять классы с исторически низким качеством распознавания (такие как \textit{bicycle}, \textit{fence}, \textit{pole} из-за их малой площади и сложной геометрии).

Интегральная метрика mIoU рассчитывается как арифметическое среднее значений IoU по всем $K$ классам:
\begin{equation}
    \text{mIoU} = \frac{1}{K} \sum_{k=1}^{K} \text{IoU}_k.
\end{equation}
В отличие от простой пиксельной точности, mIoU является робастной метрикой в условиях сильного дисбаланса классов, который типичен для дорожных сцен (где класс «дорога» может занимать значительную часть изображения, а «светофор» --- менее долей процента).

Дополнительно в работе используется базовая метрика \textbf{Pixel Accuracy (PA)} --- доля корректно классифицированных пикселей по отношению к общему числу пикселей изображения:
\begin{equation}
    \text{PA} = \frac{\sum_{k=1}^{K} TP_k}{\sum_{k=1}^{K} (TP_k + FP_k + FN_k + TN_k)} = \frac{\sum_{i=1}^{H} \sum_{j=1}^{W} \mathbb{I}(\hat{y}_{i,j} = y_{i,j})}{H \cdot W},
\end{equation}
где $\mathbb{I}(\cdot)$ --- индикаторная функция, равная $1$ при выполнении условия и $0$ в противном случае. Хотя PA легко интерпретируется, она может быть обманчиво высокой при доминировании фоновых классов, поэтому используется лишь в качестве вспомогательного индикатора.

Для оценки устойчивости моделей в \textit{adverse}-условиях (неблагоприятных условиях видимости), отдельно агрегируются и анализируются метрики по специфическим подмножествам целевого датасета (например, сценарии \textit{fog}, \textit{night}, \textit{rain}, \textit{snow} в датасете ACDC). Это позволяет явно количественно оценить степень деградации сегментации для каждого типа погодного или оптического искажения.

\subsection{Архитектурные подходы}

\subsubsection{CNN-подходы (DeepLabV3+ как классический baseline)}
Сверточные архитектуры долгое время являлись основным стандартом для задачи семантической сегментации. В работе в качестве классического подхода рассматривается \textbf{DeepLabV3+ \cite{chen2018deeplabv3plus}}, сочетающая энкодер-декодерную схему, модуль \textit{Atrous Spatial Pyramid Pooling} и \textit{dilated} свертки для расширения эффективного рецептивного поля без потери разрешения признаков. Такой подход позволяет одновременно учитывать контекст на нескольких масштабах и уточнять границы объектов, что особенно важно для дорожных сцен с большим числом тонких структур и малых объектов. В оригинальной статье авторы сообщают о результате \textbf{82.1\% mIoU} на тестовом наборе \textit{Cityscapes} и \textbf{89.0\% mIoU} на \textit{PASCAL VOC 2012} без дополнительного post-processing. DeepLabV3+ в данной работе используется как понятный и интерпретируемый CNN-baseline для сравнения с более современными трансформерными и foundation-подходами.

\subsubsection{Transformer-подходы (SegFormer)}
\textbf{SegFormer \cite{xie2021segformer} }относится к высокоэффективным трансформерным архитектурам для семантической сегментации. Особенностью архитектуры является использование иерархического энкодера, который выдает признаки на нескольких масштабах, не требуя при этом позиционного кодирования, и легкого MLP-декодера (многослойного персептрона). Декодер объединяет разномасштабные признаки для предсказания финальной маски, что обеспечивает хороший баланс между качеством сегментации и вычислительной сложностью. В оригинальной статье наиболее сильная конфигурация \textit{SegFormer-B5} достигает \textbf{84.0\% mIoU} на тестовом наборе \textit{Cityscapes}, а вариант \textit{SegFormer-B4} показывает \textbf{51.8\% mIoU} на \textit{ADE20K}. В данной работе используется конфигурация \textit{SegFormer-B2} (реализованная через библиотеку \texttt{segmentation\_models\_pytorch}), которая представляет современный практический подход и выступает в качестве сильного трансформерного baseline, не являющегося foundation-моделью в строгом смысле.

\subsubsection{CNN-based foundation backbones (InternImage)}
\textbf{InternImage \cite{wang2023internimage} }представляет собой CNN-based foundation-модель, в которой в качестве основного оператора используется \textit{DCNv3} --- деформируемая свертка, обеспечивающая большое эффективное рецептивное поле и адаптивное пространственное агрегирование признаков. Архитектура строится как иерархический четырехстадийный бэкбон, а ее блоки сочетают свойства современных CNN и transformer-подобный дизайн, включая \textit{Layer Normalization}, \textit{feed-forward network} и \textit{GELU}. Благодаря этому InternImage рассматривается как важный промежуточный вариант между классическими сверточными сетями и крупными vision foundation models.

\subsubsection{\textit{Foundation}-бэкбоны в сегментации (DINOv2)}
Foundation-подходы используют предварительное обучение на масштабных и разнообразных наборах изображений без разметки \textit{(self-supervised learning)}, что существенно повышает переносимость признаков между различными доменами и задачами. \textbf{DINOv2 \cite{oquab2023dinov2}} --- современная архитектура на базе \textit{Vision Transformer}, которая благодаря обучению с дистилляцией и специальным целевым функциям формирует богатые универсальные визуальные представления. В оригинальной статье показано, что даже замороженный бэкбон DINOv2 в связке с \textit{ViT-Adapter} достигает \textbf{60.2\% mIoU} на \textit{ADE20K}, \textbf{86.9\% mIoU} на \textit{Cityscapes} и \textbf{89.0\% mIoU} на \textit{Pascal VOC}.

В рамках исследования модель \textbf{DINOv2-Small} рассматривается как основная для проверки гипотезы о том, что предобученные универсальные признаки повышают устойчивость сегментации в условиях пониженной видимости, низкого контраста и сложных погодных условий. Для решения задачи плотного предсказания (сегментации) поверх замороженного или дообучаемого (fine-tuning) бэкбона DINOv2 применяется специально разработанный легковесный линейный декодер (\texttt{LinearDecoder}). Он преобразует одномерную последовательность патч-токенов обратно в двумерное пространственное представление и с помощью билинейной интерполяции восстанавливает признаки до исходного разрешения изображения для попиксельной классификации.

\subsection{Domain Shift: DG, UDA, FDA --- термины и различия}

\subsubsection{DG (Domain Generalization): обучение без target-домена}

\subsubsection{UDA (Unsupervised Domain Adaptation): обучение с неразмеченным target}

\subsubsection{FDA (Fourier Domain Adaptation): частотная техника адаптации}

\subsubsection{Почему в данной работе основной протокол --- source-only / DG-style}

\subsection{Обзор литературы и SOTA-контекст}

\subsubsection{Лидерборд \texorpdfstring{Cityscapes$\rightarrow$ACDC}{Cityscapes->ACDC}: что показывают UDA-методы}

\subsubsection{Какие идеи можно перенести в DG-постановку}

\subsubsection{Ограничения сопоставления DG-результатов с UDA-рейтинговыми результатами}

\section{Данные и экспериментальный протокол}

\subsection{Описание датасетов}

\subsubsection{Cityscapes: структура, классы, разметка}

\subsubsection{ACDC: adverse-условия (fog/night/rain/snow), назначение в работе}

\subsection{Протокол экспериментов}

\subsubsection{Обучение на Cityscapes}

\subsubsection{Оценка на ACDC (без обучения на ACDC в основной линии)}

\subsubsection{Риски и ограничения протокола}

\subsection{Метрики и критерии сравнения}

\subsubsection{Overall mIoU и per-condition mIoU}

\subsubsection{Class-wise IoU с акцентом на трудные/слабоконтрастные классы}

\section{Реализация системы и базовый пайплайн}

\subsection{Технический стек и инфраструктура}

\subsubsection{PyTorch / PyTorch Lightning / Albumentations}

\subsubsection{Логирование экспериментов в ClearML}

\subsection{Базовый пайплайн обучения и оценки}

\subsubsection{Подготовка данных и аугментации baseline}

\subsubsection{Конфигурации обучения и выбор checkpoint}

\subsubsection{Скрипт/ноутбук оценки на ACDC}

\subsection{Исследуемые модели}

\subsubsection{DINOv2-Small (fine-tuning)}

\subsubsection{SegFormer-B2}

\subsubsection{Роль дополнительных baseline-моделей (DeepLabV3+, InternImage) как контекст}

\section{Базовые результаты и их анализ}

\subsection{Результаты baseline на Cityscapes}

\subsubsection{Сравнение моделей по валидационным метрикам}

\subsection{Cross-domain результаты на ACDC}

\subsubsection{Overall и per-condition сравнение (fog/night/rain/snow)}

\subsubsection{Качественный анализ ошибок и типовых провалов}

\subsection{Выводы из baseline-этапа}

\subsubsection{Почему в основную линию берутся DINOv2 и SegFormer-B2}

\subsubsection{Что baseline не решает в low-contrast сценариях}

\section{План методических улучшений}

\subsection{Экспериментальная шкала E0--E4 (накопительный протокол)}

\subsubsection{E0: baseline}

\subsubsection{E1: low-contrast/weather augmentations}

\subsubsection{E2: class-balanced/focal}

\subsubsection{E3: frequency randomization (FDA-style идея без target)}

\subsubsection{E4: consistency regularization (weak/strong) на source}

\subsection{Две линии моделей в едином протоколе}

\subsubsection{DINOv2: ожидаемые эффекты и риски}

\subsubsection{SegFormer-B2: ожидаемые эффекты и риски}

\subsection{Отдельный подпункт о UDA (опционально)}

\subsubsection{Что такое UDA и как возможно обучение на неразмеченном ACDC}

\subsubsection{Почему UDA вынесен в дополнительный блок, а не в основной протокол}

\section{Заключение}

\subsection{Итоги по текущему этапу}

\subsection{Что выполнено и что планируется завершить}

\subsection{Практические рекомендации и дальнейшая работа}

\section{Приложения}

\subsection{Детальные таблицы метрик}

\subsection{Дополнительные визуализации предсказаний}

\subsection{Конфигурации экспериментов и параметры запусков}

\newpage
\printbibliography[heading=bibintoc, title={Список использованной литературы}]
\end{document}